\iffalse
\begin{abstract}
Reasoning failures in language models often arise not because training data are mislabeled, but because some examples reduce loss through shortcut updates---surface patterns, answer templates, or spurious correlations that do not support generalization. We study where this distinction appears during training and show that shortcut behavior is more directly revealed by \emph{gradient geometry} than by scalar loss: shortcut-promoting samples exhibit two observable signatures, low alignment with a small validation gradient that reflects generalizable reasoning, and concentrated gradient energy on answer tokens rather than on intermediate reasoning tokens. This yields a directional view of shortcut-aware reasoning training. The goal is not to decide which samples to keep, but which components of each sample's gradient to trust. We instantiate the view in \textbf{SART}, a detection-then-correction framework that ranks samples by a \textsc{ShortcutScore} and applies a per-sample, conflict-only projection that removes only the component opposing the validation gradient. Across three controlled reasoning benchmarks, SART improves over the strongest of eleven baselines by $+4.7$\,pp accuracy and $+7.1$\,pp robustness, with the conflict-gated projection alone contributing $+36.5$\,pp robustness in ablations shows that directional correction, not detection accuracy, is the load-bearing mechanism. A Qwen3-8B / GSM8K pilot shows the detection geometry transfers cleanly to pretrained LLMs, while correction at this scale must respect autoregressive generation constraints, which is a constrained variant of the same principle. Overall, this paper proves the principle and mechanism under controlled shortcut injection, shows that detection transfers to LLM scale, and leaves full LLM correction as a constrained scaling problem---reframing shortcut mitigation as a problem of gradient alignment and directional correction, rather than only loss shaping or sample filtering.
\end{abstract}
\fi

\begin{abstract}
Reasoning failures in language models often arise not because training data are mislabeled, but because some examples reduce loss through shortcut updates: surface patterns, answer templates, or spurious correlations that do not support generalization. We show that this failure is structurally hidden along three axes: (i) loss ambiguity, where shortcut and reasoning samples can have similar scalar losses; (ii) sample ambiguity, where one example contains both useful reasoning signal and harmful shortcut signal; and (iii) correction ambiguity, where filtering or reweighting changes which samples are used or how strongly, but not which gradient directions are trusted. We argue that reliable shortcut mitigation requires operating on gradient-level alignment rather than scalar loss or sample identity.
We propose Shortcut-Aware Reasoning Training (SART), a detection-then-correction framework with two components. First, \textsc{ShortcutScore} ranks samples by low alignment with a small validation gradient and high concentration of gradient energy on final-answer tokens. Second, SART applies a per-sample, conflict-only projection that removes only the component of a flagged gradient opposing the validation gradient, while preserving its aligned residual. Across three controlled reasoning benchmarks, SART improves over the strongest of eleven baselines by $+4.7$ pp accuracy and $+7.1$ pp robustness. Ablations show that directional correction is the load-bearing mechanism: removing conflict-only projection reduces robustness by $36.5$ pp, while scalar reweighting achieves higher detection F1 but worse generalization. A Qwen3-8B/GSM8K pilot shows that detection transfers to pretrained LLMs, while full correction must respect autoregressive generation constraints. These results suggest that shortcut mitigation does not lack useful training signal; it lacks directional discrimination.
\end{abstract}
